\documentclass[11pt]{article}

\usepackage{amsmath, amssymb}
\usepackage{geometry}
\usepackage{setspace}
\usepackage{titlesec}
\usepackage{hyperref}

\geometry{margin=1in}
\setstretch{1.1}

\titleformat{\section}
  {\normalfont\Large\bfseries}{\thesection}{1em}{}

\titleformat{\subsection}
  {\normalfont\large\bfseries}{\thesubsection}{1em}{}

\begin{document}

%%%%%%%%%%%%%%%%%%%%%%%%%%%%%%%%%%%%%%%%%%%%%%%%%%%%
% Header
%%%%%%%%%%%%%%%%%%%%%%%%%%%%%%%%%%%%%%%%%%%%%%%%%%%%

\begin{center}
{\large \textbf{CSE 400 : Fundamentals of Probability in Computing}}\\
\vspace{0.3cm}
Nidhi Mehta – AU2440019\\
January 29, 2026
\end{center}

\vspace{0.5cm}

\begin{center}
{\LARGE \textbf{Lecture Scribe}}\\[0.2cm]
{\large Lecture 8: Gaussian, Uniform, Exponential, and Gamma Random Variables}
\end{center}

\vspace{0.8cm}

%%%%%%%%%%%%%%%%%%%%%%%%%%%%%%%%%%%%%%%%%%%%%%%%%%%%
\section*{Gaussian Random Variable}

%%%%%%%%%%%%%%%%%%%%%%%%%%%%%%%%%%%%%%%%%%%%%%%%%%%%
\subsection*{1. Definition and Properties}

\textbf{Definition 1.1 (Gaussian Random Variable)}

A Gaussian random variable is one whose PDF can be written as:

\[
f_X(x) =
\frac{1}{\sqrt{2\pi\sigma^2}}
\exp\left(
-\frac{(x-m)^2}{2\sigma^2}
\right)
\]

where:

\begin{itemize}
    \item $m$ = mean
    \item $\sigma^2$ = variance
    \item $\sigma$ = standard deviation
\end{itemize}

We denote:

\[
X \sim \mathcal{N}(m,\sigma^2)
\]

\subsection*{Interpretation}

\begin{itemize}
    \item Bell-shaped curve
    \item Symmetric about mean $m$
    \item Total area under curve = 1
    \item CDF is S-shaped and monotonically increasing
\end{itemize}

\subsection*{Properties}

1. Symmetry:
\[
f_X(m+a) = f_X(m-a)
\]

2. Mean:
\[
E[X] = m
\]

3. Variance:
\[
\text{Var}(X) = \sigma^2
\]

%%%%%%%%%%%%%%%%%%%%%%%%%%%%%%%%%%%%%%%%%%%%%%%%%%%%
\newpage
%%%%%%%%%%%%%%%%%%%%%%%%%%%%%%%%%%%%%%%%%%%%%%%%%%%%

\section*{Standard Forms}

\subsection*{Error Function}

\[
\text{erf}(x) =
\frac{2}{\sqrt{\pi}}
\int_0^x e^{-t^2} dt
\]

\subsection*{Complementary Error Function}

\[
\text{erfc}(x)
=
1 - \text{erf}(x)
=
\frac{2}{\sqrt{\pi}}
\int_x^\infty e^{-t^2} dt
\]

\subsection*{$\Phi$-Function (CDF of Standard Normal)}

\[
\Phi(x)
=
\frac{1}{\sqrt{2\pi}}
\int_{-\infty}^{x}
e^{-t^2/2} dt
\]

\subsection*{$Q$-Function (Gaussian Tail Function)}

\[
Q(x)
=
\frac{1}{\sqrt{2\pi}}
\int_x^\infty
e^{-t^2/2} dt
\]

\subsection*{Key Identity}

\[
Q(x) = 1 - \Phi(x)
\]

%%%%%%%%%%%%%%%%%%%%%%%%%%%%%%%%%%%%%%%%%%%%%%%%%%%%
\newpage
%%%%%%%%%%%%%%%%%%%%%%%%%%%%%%%%%%%%%%%%%%%%%%%%%%%%

\section*{Relation Between $\Phi$ and $Q$}

\subsection*{Evaluating the CDF}

\[
F_X(x)
=
\int_{-\infty}^{x}
\frac{1}{\sqrt{2\pi\sigma^2}}
\exp
\left(
-\frac{(y-m)^2}{2\sigma^2}
\right)
dy
\]

Substitute:

\[
t = \frac{y-m}{\sigma}
\]

Then:

\[
F_X(x)
=
\Phi
\left(
\frac{x-m}{\sigma}
\right)
\]

\subsection*{Tail Probability}

\[
\Pr(X > x)
=
Q
\left(
\frac{x-m}{\sigma}
\right)
\]

\subsection*{Key Identities}

\[
Q(x) = 1 - \Phi(x)
\]

\[
F_X(x)
=
1 -
Q
\left(
\frac{x-m}{\sigma}
\right)
\]

%%%%%%%%%%%%%%%%%%%%%%%%%%%%%%%%%%%%%%%%%%%%%%%%%%%%
\newpage
%%%%%%%%%%%%%%%%%%%%%%%%%%%%%%%%%%%%%%%%%%%%%%%%%%%%

\section*{Worked Example}

Given:

\[
f_X(x)
=
\frac{1}{\sqrt{8\pi}}
\exp
\left(
-\frac{(x+3)^2}{8}
\right)
\]

\subsection*{Step 1: Identify Mean and Variance}

Compare with:

\[
f_X(x)=
\frac{1}{\sqrt{2\pi\sigma^2}}
\exp
\left(
-\frac{(x-m)^2}{2\sigma^2}
\right)
\]

We identify:

\[
m = -3
\]

\[
2\sigma^2 = 8
\]

\[
\sigma^2 = 4
\]

\[
\sigma = 2
\]

\subsection*{1) $\Pr(X \le 0)$}

\[
\Pr(X \le 0)
=
F_X(0)
=
\Phi
\left(
\frac{0-(-3)}{2}
\right)
=
\Phi
\left(
\frac{3}{2}
\right)
=
1 -
Q
\left(
\frac{3}{2}
\right)
\]

\subsection*{2) $\Pr(X > 4)$}

\[
\Pr(X > 4)
=
Q
\left(
\frac{4-(-3)}{2}
\right)
=
Q
\left(
\frac{7}{2}
\right)
\]

(Other parts follow same standardization method.)

%%%%%%%%%%%%%%%%%%%%%%%%%%%%%%%%%%%%%%%%%%%%%%%%%%%%
\newpage
%%%%%%%%%%%%%%%%%%%%%%%%%%%%%%%%%%%%%%%%%%%%%%%%%%%%

\section*{Summary Section}

\subsection*{Key Flashcard Concepts}

\begin{itemize}
    \item Gaussian $\rightarrow$ Bell shape
    \item Mean $\rightarrow$ Center
    \item Variance $\rightarrow$ Spread
    \item $\Phi$ $\rightarrow$ Left tail
    \item $Q$ $\rightarrow$ Right tail
    \item Standardization $\rightarrow$ Subtract mean, divide by $\sigma$
\end{itemize}

\subsection*{Core Formulas}

\[
f_X(x)=
\frac{1}{\sqrt{2\pi\sigma^2}}
e^{-(x-m)^2/2\sigma^2}
\]

\[
F_X(x)=
\Phi
\left(
\frac{x-m}{\sigma}
\right)
\]

\[
\Pr(X>x)=
Q
\left(
\frac{x-m}{\sigma}
\right)
\]

\[
Q(x)=1-\Phi(x)
\]

\subsection*{What To Absorb}

\begin{itemize}
    \item Always standardize before applying tables
    \item Use $\Phi$ for left probability
    \item Use $Q$ for right probability
    \item Identify mean and variance correctly
\end{itemize}

\subsection*{Applications in Exams}

\begin{itemize}
    \item Probability evaluation
    \item Noise modeling
    \item Signal detection
    \item Statistical modeling
\end{itemize}

\subsection*{Tricks}

\begin{itemize}
    \item “Greater than” $\rightarrow$ Use $Q$
    \item “Less than” $\rightarrow$ Use $\Phi$
    \item If PDF has $(x+a)^2$ $\rightarrow$ Mean is $-a$
\end{itemize}

\subsection*{Be Careful}

\begin{itemize}
    \item Do not confuse $\sigma$ and $\sigma^2$
    \item Don’t forget negative sign inside square
    \item Always standardize properly
    \item $Q(x) \neq \Phi(x)$
\end{itemize}

%%%%%%%%%%%%%%%%%%%%%%%%%%%%%%%%%%%%%%%%%%%%%%%%%%%%
\end{document}
